\documentclass[12pt, letterpaper]{article}

% ─── Paquetes ────────────────────────────────────────────────────────────────
\usepackage[utf8]{inputenc}
\usepackage[T1]{fontenc}
\usepackage[spanish]{babel}
\usepackage{geometry}
  \geometry{margin=2.5cm}
\usepackage{booktabs}
\usepackage{longtable}
\usepackage{graphicx}
  \graphicspath{{../03_figuras/}}
\usepackage{amsmath, amssymb}
\usepackage{siunitx}
  \sisetup{locale=ES, output-decimal-marker={,}}
\usepackage{hyperref}
  \hypersetup{colorlinks=true, linkcolor=black, urlcolor=blue, citecolor=black}
\usepackage{listings}
\usepackage{xcolor}
\usepackage{fancyhdr}
\usepackage{caption}
\usepackage{subcaption}
\usepackage{float}

% ─── Estilo de código Python ──────────────────────────────────────────────────
\definecolor{codeblue}{rgb}{0.13,0.29,0.53}
\definecolor{codegray}{rgb}{0.5,0.5,0.5}
\definecolor{codegreen}{rgb}{0.0,0.5,0.0}
\lstdefinestyle{python}{
  language=Python,
  basicstyle=\ttfamily\footnotesize,
  keywordstyle=\color{codeblue}\bfseries,
  commentstyle=\color{codegray}\itshape,
  stringstyle=\color{codegreen},
  numbers=left, numberstyle=\tiny\color{codegray},
  breaklines=true, frame=single,
  showstringspaces=false,
}

% ─── Encabezado y pie de página ───────────────────────────────────────────────
\pagestyle{fancy}
\fancyhf{}
\lhead{\small WASAFF Consulting — Informe Técnico Confidencial}
\rhead{\small Kaeser / CMPC Puente Alto}
\cfoot{\thepage}

% ─── Metadatos ────────────────────────────────────────────────────────────────
\title{
  \vspace{-1cm}
  {\large \textbf{WASAFF Consulting}} \\[0.3cm]
  {\LARGE Modelo Matemático: Sistema de Recuperación de Calor} \\[0.2cm]
  {\large Sala de Compresores Kaeser --- Papeles Cordillera (CMPC Puente Alto)}
}
\author{
  Felipe Wasaff, Director Técnico \\
  \small \texttt{felipe@wasaffconsulting.com}
}
\date{Abril 2026}

% ─────────────────────────────────────────────────────────────────────────────
\begin{document}

\maketitle
\thispagestyle{empty}

\vspace{0.5cm}
\begin{center}
  \small
  \textbf{Confidencial} --- Este documento es de uso exclusivo de
  Kaeser Compresores de Chile SpA y Papeles Cordillera S.A. (CMPC).
  Prohibida su reproducción o distribución sin autorización escrita de WASAFF Consulting.
\end{center}

\tableofcontents
\newpage

% ─────────────────────────────────────────────────────────────────────────────
\section{Resumen Ejecutivo}
% Máximo 1 página. Escribir al final, cuando todos los resultados estén listos.
% Incluir: problema planteado, metodología en 2 líneas, resultado principal,
%          recomendación clave, valor de la recuperación en kW y en \$/año estimado.

\textit{[Completar al cierre del proyecto]}

% ─────────────────────────────────────────────────────────────────────────────
\section{Antecedentes y Datos del Sistema}

\subsection{Descripción del sistema}

La planta Papeles Cordillera (CMPC Puente Alto) opera una sala de compresores
Kaeser de tornillo rotatorio con capacidad de recuperar calor residual para
transferirlo al circuito de agua industrial. Los dos circuitos deben permanecer
hidráulicamente aislados.

\subsection{Compresores instalados}

\begin{table}[H]
  \centering
  \caption{Parámetros térmicos e hidráulicos de compresores (fuente: datasheets Kaeser)}
  \label{tab:compresores}
  \begin{tabular}{lccccc}
    \toprule
    Modelo & Cant. & $P_\text{nom}$ (kW) & $\dot{Q}$ (kW) &
    $\dot{V}$ a $\Delta T=\SI{25}{\kelvin}$ (m³/h) & $\Delta P$ (bar) \\
    \midrule
    FSD575  & 3 & 315 & 246 & 8{,}61 & 0{,}40 \\
    DSDX305 & 2 & 160 & 130 & 4{,}50 & 0{,}25 \\
    ESD445  & 1 & 250 & 196 & 6{,}71 & 0{,}40 \\
    \bottomrule
  \end{tabular}
\end{table}

\noindent\textbf{Nota operacional:} El compresor FSD575 N°3 opera como 100\% de respaldo;
no funciona simultáneamente con FSD575 N°1 o N°2.

\subsection{Escenarios de operación}

\begin{table}[H]
  \centering
  \caption{Escenarios de operación del sistema}
  \label{tab:escenarios}
  \begin{tabular}{lccc}
    \toprule
    Escenario & Compresores activos & $\dot{Q}_\text{total}$ (kW) & $\dot{V}_\text{total}$ (m³/h) \\
    \midrule
    Normal (80\% del tiempo) & FSD575×2 + DSDX305 (1-2-4) & 622 & 21{,}72 \\
    Mínimo & \textit{por confirmar} & --- & --- \\
    Máximo & \textit{por confirmar} & --- & --- \\
    \bottomrule
  \end{tabular}
\end{table}

\subsection{Agua industrial CMPC (circuito secundario)}

\begin{table}[H]
  \centering
  \caption{Parámetros del circuito de agua industrial CMPC}
  \label{tab:agua_cmpc}
  \begin{tabular}{lcc}
    \toprule
    Parámetro & Mínimo & Máximo \\
    \midrule
    Temperatura (\si{\degreeCelsius}) & 6 & 16 \\
    Presión (bar) & 3{,}7 & 4{,}2 \\
    Caudal (m³/h) & 27 & 50 \\
    \bottomrule
  \end{tabular}
\end{table}

% ─────────────────────────────────────────────────────────────────────────────
\section{Alcance y Supuestos}

\subsection{Lo que se modela}
\begin{itemize}
  \item Análisis $\varepsilon$-NTU del intercambiador de aislamiento hidráulico
  \item Selección y dimensionamiento del intercambiador de placas
  \item Red hidráulica completa (Darcy-Weisbach + Colebrook-White iterativo)
  \item Dimensionamiento de bomba de circulación
  \item Dimensionamiento de acumulador hidráulico
  \item Verificación de tres escenarios: mínimo, normal y máximo
\end{itemize}

\subsection{Fuera de alcance}
\begin{itemize}
  \item Diseño estructural de soportes y obra civil
  \item Selección de proveedor específico de equipos
  \item Especificación de instrumentación y control
  \item Visita a terreno (cotizable por separado)
\end{itemize}

\subsection{Supuestos principales}
\begin{itemize}
  \item Fluido de trabajo: agua líquida en ambos circuitos
  \item Propiedades termofísicas del agua según correlaciones ASHRAE 2021
        (rango válido \SI{0}{\degreeCelsius}--\SI{100}{\degreeCelsius})
  \item Flujo en régimen permanente para cada escenario de operación
  \item Intercambiador modelado como contracorriente puro
  \item Pérdidas de calor al ambiente despreciables (sistema aislado)
  \item \textit{[TODO: Felipe --- agregar supuestos de geometría de tuberías
        una vez disponible el layout]}
\end{itemize}

% ─────────────────────────────────────────────────────────────────────────────
\section{Metodología}

\subsection{Análisis \texorpdfstring{$\varepsilon$}{ε}-NTU del intercambiador}

El método $\varepsilon$-NTU relaciona la efectividad del intercambiador con el
número de unidades de transferencia sin requerir conocer las temperaturas de salida
de antemano \cite{incropera2011}.

\begin{equation}
  \varepsilon = \frac{\dot{Q}}{\dot{Q}_\text{max}}
  \qquad \text{donde} \qquad
  \dot{Q}_\text{max} = C_\text{min}(T_{h,\text{in}} - T_{c,\text{in}})
  \label{eq:efectividad}
\end{equation}

Para un intercambiador a contracorriente con $C_r = C_\text{min}/C_\text{max} \neq 1$:

\begin{equation}
  \text{NTU} = \frac{\ln\!\left[\dfrac{\varepsilon - 1}{\varepsilon C_r - 1}\right]}{C_r - 1}
  \label{eq:NTU_contracorriente}
\end{equation}

El producto $UA$ se obtiene como:
\begin{equation}
  UA = \text{NTU} \cdot C_\text{min}
  \label{eq:UA}
\end{equation}

\subsection{Red hidráulica: Darcy-Weisbach + Colebrook-White}

La pérdida de carga en cada tramo se calcula con:

\begin{equation}
  \Delta P = f \cdot \frac{L}{D} \cdot \frac{\rho v^2}{2}
  \label{eq:darcy_weisbach}
\end{equation}

donde el factor de fricción $f$ en régimen turbulento se obtiene resolviendo
iterativamente la ecuación de Colebrook-White \cite{white2011}:

\begin{equation}
  \frac{1}{\sqrt{f}} = -2\log_{10}\!\left(\frac{\varepsilon_r}{3{,}7\,D} +
  \frac{2{,}51}{Re\sqrt{f}}\right)
  \label{eq:colebrook}
\end{equation}

La iteración se resuelve por Newton-Raphson con criterio de convergencia
$|f_{k+1} - f_k| < 10^{-8}$.

% ─────────────────────────────────────────────────────────────────────────────
\section{Resultados}
% Completar con tablas y figuras generadas por resultados.py

\subsection{Intercambiador de aislamiento}

\textit{[Tabla de resultados ε-NTU por escenario --- generada por \texttt{resultados.py}]}

\subsection{Red hidráulica}

\textit{[Tabla de pérdidas por tramo y ΔP total --- generada por \texttt{resultados.py}]}

\subsection{Bomba de circulación}

\textit{[Punto de operación requerido y recomendación de configuración]}

\subsection{Acumulador}

\textit{[Volumen y presiones de diseño]}

% ─────────────────────────────────────────────────────────────────────────────
\section{Análisis y Discusión}

\textit{[Completar al cierre del proyecto]}

% ─────────────────────────────────────────────────────────────────────────────
\section{Conclusiones y Recomendaciones}

\textit{[Completar al cierre del proyecto]}

% ─────────────────────────────────────────────────────────────────────────────
\begin{thebibliography}{9}

\bibitem{incropera2011}
  Incropera, F.P., DeWitt, D.P., Bergman, T.L., Lavine, A.S.
  \textit{Fundamentals of Heat and Mass Transfer}, 7\textsuperscript{a} ed.
  John Wiley \& Sons, 2011.

\bibitem{white2011}
  White, F.M.
  \textit{Fluid Mechanics}, 8\textsuperscript{a} ed.
  McGraw-Hill, 2011.

\bibitem{ashrae2021}
  ASHRAE.
  \textit{ASHRAE Handbook --- Fundamentals}.
  American Society of Heating, Refrigerating and
  Air-Conditioning Engineers, 2021.

\end{thebibliography}

% ─────────────────────────────────────────────────────────────────────────────
\appendix

\section{Código Python --- \texttt{main.py}}
\lstinputlisting[style=python, caption={main.py --- punto de entrada}]
  {../02_analisis/main.py}

\section{Memoria de cálculo --- Red hidráulica}
% Desarrollar aquí los cálculos de un tramo representativo a mano
% para validar el código.

\end{document}
